% Closure of wp3 lem:remainder  [drafted v2 — author line-check required]
\closuredraft
\begin{proof}
Fix a compact $K\subset\R^2$ and let $R_K:=\sup_{h\in K}\norm h$. For $N,n$
large, $\vzero+\Dn h$ lies in the interior of $\Theta$ for all $h\in K$
because $\Dn\to0$. A third-order Taylor expansion of
$h\mapsto\elln(\vzero+\Dn h)$ around $h=0$ gives, for every $h\in K$,
\begin{equation}\label{eq:rem-taylor}
 \elln(\vzero+\Dn h)-\elln(\vzero)
 =h^\top\Delta_{N,n}
 +\tfrac12 h^\top\big[\Dn\nabla^2_\vartheta\elln(\vzero)\Dn\big]h
 +\tfrac16\,\mathrm T_{N,n}(\bar h)[h,h,h],
\end{equation}
where $\bar h$ lies on the segment $[0,h]$ and $\mathrm T_{N,n}(\bar h)$ is
the trilinear form with entries
$\Dn^{(a)}\Dn^{(b)}\Dn^{(c)}\,
\partial^3_{abc}\elln(\vzero+\Dn\bar h)$, $a,b,c\in\{\theta,\sigma\}$,
$\Dn^{(\theta)}=\LamN^{-1/2}$, $\Dn^{(\sigma)}=(Nn)^{-1/2}$.
By \eqref{eq:observed-info} (Step 4 of the proof of Lemma~\ref{lem:score-clt}),
the quadratic term is $-\frac12h^\top J(\vzero)h+o_{\Pp}(1)$ uniformly on $K$
(a fixed quadratic form applied to a compact set). It therefore suffices to
show
\begin{equation}\label{eq:rem-goal}
 \sup_{\norm{\bar h}\le R_K}\;
 \max_{a,b,c}\;
 \Big|\Dn^{(a)}\Dn^{(b)}\Dn^{(c)}\,
 \partial^3_{abc}\elln(\vzero+\Dn\bar h)\Big|
 =o_{\Pp}(1).
\end{equation}

\emph{Step 1: structure of third derivatives.}
Per transition the log-density is
$-\frac12\log v(\vartheta)-r(\theta)^2/(2v(\vartheta))+\text{const}$ with
$r(\theta)=Y_{k,j+1}-\rho_{k,n}(\theta)Y_{k,j}$ and
$v=\sigma^2q_k^2h_{k,n}(\theta)$. All $\theta$-derivatives of $r$ obey
$|\partial_\theta^m r|=(\lambda_k\Delta_n)^m\rho_{k,n}(\theta)|Y_{k,j}|$, and,
uniformly on $\Theta$ for $\lambda_k\Delta_n\le x_0$,
\begin{equation}\label{eq:rem-logv}
 |\partial_\theta^m\log v|\le C(\lambda_k\Delta_n)^m\ \ (m=1,2,3),
 \qquad
 |\partial_\sigma\log v|=2/\sigma,\quad
 \partial_\sigma^2\log v=-2/\sigma^2,\ \ldots
\end{equation}
(the first bound extends \eqref{eq:g-exp} to two more derivatives of
$x\mapsto\log\{(1-e^{-2\theta x})/(2\theta x)\}$, which is analytic with all
derivatives $O(x^m)$ at $0$, uniformly for $\theta$ compact). Expanding
$\partial^3_{abc}$ of the per-transition density by the chain rule, every
resulting term is a product of factors from \eqref{eq:rem-logv} and of
$r^2/v$, $\partial_\theta r\cdot r/v$, $(\partial_\theta r)^2/v$,
$\partial_\theta^2r\cdot r/v$, etc. Collecting powers, one obtains the
per-transition envelope, valid uniformly on $\Theta$:
\begin{equation}\label{eq:rem-envelope}
 |\partial^3_{abc}\ell_{k,j}(\vartheta)|
 \le C\,(\lambda_k\Delta_n)^{m_\theta(a,b,c)}\,\sigma^{-m_\sigma(a,b,c)}
 \Big(1+\frac{r(\theta)^2}{v(\vartheta)}
 +\frac{\lambda_k\Delta_n\,Y_{k,j}^2}{v(\vartheta)/\Delta_n}\Big),
\end{equation}
where $m_\theta$ counts $\theta$-indices and $m_\sigma$ counts
$\sigma$-indices among $\{a,b,c\}$. (For the mean part, each $\theta$-index
contributes one factor $\lambda_k\Delta_n$; terms with
$(\partial_\theta r)^2/v=\lambda_k^2\Delta_n^2\rho^2Y^2/v$ are absorbed in the
last bracket using $v\ge c\sigma^2q_k^2\Delta_n$ and
$Y^2/(q_k^2)\cdot\lambda_k\Delta_n\le$ the stated form.)

\emph{Step 2: uniformity over the shrinking box.}
For $\vartheta=\vzero+\Dn\bar h$, $\norm{\bar h}\le R_K$:
$|\theta-\theta_0|\le R_K\LamN^{-1/2}$ and
$|\sigma-\sigma_0|\le R_K(Nn)^{-1/2}$, so
$\sigma^{-1}\le C$, and for every $k\le N$,
\[
 \frac{v(\vzero)}{v(\vartheta)}\le C,
 \qquad
 \frac{r(\theta)^2}{v(\vartheta)}
 \le C\Big(\varepsilon_{k,j}^2
 +\frac{(\rho_{k,n}(\theta_0)-\rho_{k,n}(\theta))^2Y_{k,j}^2}{v(\vzero)}\Big)
 \le C\Big(\varepsilon_{k,j}^2
 +\frac{\lambda_k^2\Delta_n^2Y_{k,j}^2}{\LamN\,v(\vzero)}R_K^2\Big),
\]
using the mean value theorem for $\rho_{k,n}$ and
$e^{\pm\lambda_k\Delta_nR_K\LamN^{-1/2}}\le C$. Hence the envelope
\eqref{eq:rem-envelope} holds with $\vartheta$-free right-hand side
\begin{equation}\label{eq:rem-envelope2}
 E_{k,j}
 :=C(\lambda_k\Delta_n)^{m_\theta}
 \Big(1+\varepsilon_{k,j}^2
 +\frac{\lambda_k\Delta_n\,Y_{k,j}^2}{\sigma_0^2q_k^2}\Big),
 \qquad
 \E E_{k,j}\le C(\lambda_k\Delta_n)^{m_\theta}
 \Big(1+\frac{u_{0,k}^2\rho_k^{2j}}{s_k^2}\Big)
\end{equation}
(using $\E Y_{k,j}^2\le u_{0,k}^2\rho_k^{2j}+s_k^2$ and
$\lambda_k\Delta_n s_k^2/(\sigma_0^2q_k^2)=\Delta_n/(2\theta_0)\cdot
\lambda_k\le C$ after multiplying by $\lambda_k\Delta_n\le x_0$; the
initial-condition part sums over $j$ to
$Cu_{0,k}^2/(s_k^2\lambda_k\Delta_n)
=C\lambda_ku_{0,k}^2/(q_k^2)\cdot(\sigma_0^2\lambda_k\Delta_n)^{-1}\cdot
2\theta_0\lambda_k\Delta_n\cdot$(const)$\le C\lambda_ku_{0,k}^2/q_k^2$).

\emph{Step 3: rate bookkeeping.}
Summing \eqref{eq:rem-envelope2} over $k\le N$, $j<n$ and applying the scale
factors, Markov's inequality reduces \eqref{eq:rem-goal} to the vanishing of
the deterministic quantities
\[
 D_{abc}:=\Dn^{(a)}\Dn^{(b)}\Dn^{(c)}
 \Big\{n\Delta_n^{m_\theta}\sum_{k\le N}\lambda_k^{m_\theta}
 +\Delta_n^{m_\theta}\sum_{k\le N}\lambda_k^{m_\theta}
 \cdot\frac{\lambda_ku_{0,k}^2}{q_k^2}\Big\}.
\]
Using $\sum_k\lambda_k^m\le\lambda_N^{m-1}\LamN$ and
$R_N=\sum_k\lambda_ku_{0,k}^2/q_k^2=o(\LamN)$:
\begin{align*}
 (\theta\theta\theta):\quad
 &\LamN^{-3/2}\,T\Delta_n^2{\textstyle\sum}\lambda_k^3
 \le T(\lambda_N\Delta_n)^2\LamN^{-1/2}\to0;\\
 (\theta\theta\sigma):\quad
 &\LamN^{-1}(Nn)^{-1/2}\,T\Delta_n{\textstyle\sum}\lambda_k^2
 \le T\lambda_N\Delta_n\,(Nn)^{-1/2}\to0;\\
 (\theta\sigma\sigma):\quad
 &\LamN^{-1/2}(Nn)^{-1}\,T{\textstyle\sum}\lambda_k
 =T\sqrt{\LamN/(Nn)}\,(Nn)^{-1/2}
 \le T\sqrt{\lambda_N\Delta_n/T}\,(Nn)^{-1/2}\to0;\\
 (\sigma\sigma\sigma):\quad
 &(Nn)^{-3/2}\,CNn=C(Nn)^{-1/2}\to0,
\end{align*}
and the $u_0$-terms are smaller by the same comparisons with
$R_N=o(\LamN)$ in place of one factor $T\LamN/\Delta_n^{-1}$
(e.g.\ for $(\theta\theta\theta)$:
$\LamN^{-3/2}\Delta_n^3\lambda_N^2R_N\cdot\Delta_n^{-1}\cdot$(const)
$\le(\lambda_N\Delta_n)^2\LamN^{-3/2}R_N\to0$).
Since the bounds are monotone in $\norm{\bar h}\le R_K$ and polynomial in
$R_K$, the convergence is uniform on $K$; \eqref{eq:remainder} follows from
\eqref{eq:rem-taylor}.
\end{proof}
